% ============================================================
%  The Particle Model — Lesson Booklet (one section per lesson: A Reading · B Theory · C Practice · D Above level)
%  Year 7 (Year 8 extension)  |  Victorian Curriculum F–10 Version 2.0
%  Compile with: xelatex three times, then python3 fill.py Particle_Model_Lesson_Booklet_Yr7
% ============================================================
\documentclass[11pt,a5paper,twoside]{article}
\newcommand{\wrcoverfile}{../../../covers/water-resources.jpg}\newcommand{\unitname}{Water Cycle and Resources}\newcommand{\unitcaps}{WATER CYCLE AND RESOURCES}

% ── Fonts (compile with LuaLaTeX) ────────────────────────────────────────────
\usepackage{fontspec}
% Bundled fonts (in ./fonts/) so the document compiles anywhere with XeLaTeX/LuaLaTeX
\setmainfont{Lato}[                   % body — humanist sans (stands in for deck's Inter)
  Path=../fonts/, Extension=.ttf, Ligatures=TeX,
  UprightFont=Lato-Regular, ItalicFont=Lato-Italic,
  BoldFont=Lato-Bold, BoldItalicFont=Lato-BoldItalic ]
\newfontfamily\headingfont{Poppins}[  % display/headings — matches the slide decks
  Path=../fonts/, Extension=.ttf, Ligatures=TeX,
  UprightFont=Poppins-Regular, ItalicFont=Poppins-Italic,
  BoldFont=Poppins-Bold, BoldItalicFont=Poppins-BoldItalic ]
\setlength{\headheight}{22pt}

\usepackage[protrusion=true,expansion=false]{microtype}
\usepackage[a5paper, inner=17mm, outer=13mm, top=19mm, bottom=17mm, headheight=14pt, headsep=4mm, footskip=8mm]{geometry}
\usepackage{xcolor}
\usepackage{tikz}
\usetikzlibrary{arrows.meta, shapes.geometric, calc, positioning, decorations.pathreplacing}
\usepackage{tcolorbox}
\tcbuselibrary{skins, breakable, raster}
\usepackage{fancyhdr}
\usepackage{titlesec}
\usepackage{enumitem}
\usepackage{booktabs}
\usepackage{tabularx}
\usepackage{array}
\usepackage{colortbl}
\usepackage{multicol}
\usepackage{amsmath}
\usepackage{needspace}
% TikZ diagram macros are embedded below (just before \begin{document}) for easy editing.
\usepackage{longtable}
\usepackage{graphicx}
\usepackage{pifont}
\usepackage{lastpage}
\usepackage{calc}
\usepackage{mdframed}
\usepackage{setspace}
\usepackage{ragged2e}
% PDF bookmarks: one per lesson (made in \sectiondiv) plus the cover and
% contents, and links from the contents table to each lesson.
% implicit=false: hyperref adds no anchors of its own (an anchor on each
% list item nudged lists down 1.5 mm, which would undo the page fill).
% hidelinks keeps the links invisible, so the printed pages are unchanged.
\usepackage[implicit=false, hidelinks,
  pdftitle={Water Cycle and Resources — Lesson Booklet},
  pdfsubject={Years 7–8 Science, Victorian Curriculum F–10 v2.0}]{hyperref}
\usepackage{bookmark} % writes the outline itself, which implicit=false switches off in hyperref
% hyperref 7.01 with implicit=false leaves \Hy@PageAnchorSlide undefined, which
% throws at \end{document} and stops fill.py (it builds with -halt-on-error).
\makeatletter\providecommand\Hy@PageAnchorSlide{}\makeatother

% ── Colour Palette ───────────────────────────────────────────────────────────
% Remapped to the slide-deck palette (names kept so body + diagrams cascade)
\definecolor{NavyBlue}   {HTML}{16305C}
\definecolor{MidBlue}    {HTML}{2E6FE0}
\definecolor{LightBlue}  {HTML}{E8F1FF}
\definecolor{TealGreen}  {HTML}{0F6E63}
\definecolor{LightGreen} {HTML}{E0F7F3}
\definecolor{GoldDark}   {HTML}{C77F12}
\definecolor{GoldLight}  {HTML}{FFF3DC}
\definecolor{PurpleDark} {HTML}{5B3FC0}
\definecolor{PurpleLight}{HTML}{EFE9FF}
\definecolor{OrangeDark} {HTML}{C9483D}
\definecolor{OrangeLight}{HTML}{FFE9E6}
\definecolor{MidGrey}    {HTML}{AEBED4}
\definecolor{LightGrey}  {HTML}{F2F5FA}
\definecolor{DarkGrey}   {HTML}{3A4759}
\definecolor{TableHead}  {HTML}{16305C}
\definecolor{RowAlt}     {HTML}{EEF4FF}
% Bright accent tints (deck) for bars, rules and chips
\definecolor{Blue}   {HTML}{2E6FE0}
\definecolor{Teal}   {HTML}{18A999}
\definecolor{Amber}  {HTML}{F5A623}
\definecolor{Coral}  {HTML}{FF6B5E}
\definecolor{Violet} {HTML}{7C5CE0}
\definecolor{Sky}    {HTML}{4DA3FF}
\definecolor{Ink}    {HTML}{16223A}
\definecolor{Line}   {HTML}{DDE5F0}
\definecolor{Paper}  {HTML}{FBFCFE}

% ── Page Style ───────────────────────────────────────────────────────────────
\pagestyle{fancy}
\fancyhf{}
\renewcommand{\headrulewidth}{0.4pt}
\renewcommand{\headrule}{\hbox to\headwidth{\color{NavyBlue}\leaders\hrule height \headrulewidth\hfill}}
\fancyhead[L]{\headingfont\small\color{NavyBlue}The Particle Model}
\fancyhead[R]{\headingfont\small\color{DarkGrey}Year 7 \textbar{} Victorian Curriculum F--10 v2.0}
\fancyfoot[C]{\headingfont\small\color{DarkGrey}\thepage\ \textbar{} \pageref{LastPage}}
\fancypagestyle{plain}{%
  \fancyhf{}
  \fancyfoot[C]{\headingfont\small\color{DarkGrey}\thepage\ \textbar{} \pageref{LastPage}}
  \renewcommand{\headrulewidth}{0pt}
}

% ── Section Formatting ───────────────────────────────────────────────────────
\titleformat{\section}
  {\headingfont\color{NavyBlue}\Large\bfseries}
  {}{0em}{}
  [{\color{Teal}\rule[0.6ex]{2.2em}{3pt}}{\color{Line}\rule[0.6ex]{\dimexpr\linewidth-2.2em\relax}{1pt}}]
\titlespacing{\section}{0pt}{18pt}{8pt}

\titleformat{\subsection}
  {\headingfont\color{MidBlue}\large\bfseries}
  {\thesubsection}{0.6em}{}
\titlespacing{\subsection}{0pt}{12pt}{4pt}

\titleformat{\subsubsection}
  {\headingfont\color{NavyBlue}\normalsize\bfseries}
  {}{0em}{}
\titlespacing{\subsubsection}{0pt}{8pt}{2pt}

% Deck-style eyebrow / kicker (optional use)
\newcommand{\kicker}[2][Teal]{%
  {\headingfont\bfseries\footnotesize\color{#1}%
   \textcolor{#1}{\rule[0.18ex]{1.6em}{2.2pt}}\hspace{0.5em}\MakeUppercase{#2}\par}%
  \vspace{2pt}}

% ── Custom tcolorbox Environments ────────────────────────────────────────────

% Section divider banner
\newtcolorbox{sectionbanner}[2]{
  enhanced, breakable=false,
  colback=NavyBlue, colframe=NavyBlue,
  fontupper=\color{white}\bfseries\large,
  left=4mm, right=4mm, top=3mm, bottom=3mm,
  title={\color{white}\Huge\bfseries #1},
  attach boxed title to top left={yshift=-2mm, xshift=4mm},
  boxed title style={colback=NavyBlue, colframe=NavyBlue},
  after upper={\par\smallskip\normalfont\normalsize\color{white!80!NavyBlue}\textit{#2}}
}

% Key info box (blue)
\newtcolorbox{infobox}[1]{
  breakable, arc=2.5mm, boxrule=0.5pt,
  colback=LightBlue, colframe=Blue!40,
  fonttitle=\headingfont\bfseries, coltitle=MidBlue,
  title={#1},
  left=4mm, right=3.5mm, top=2.5mm, bottom=2.5mm
}

% Scaffold box (yellow)
\newtcolorbox{scaffold}[1]{
  breakable, arc=2.5mm, boxrule=0.5pt,
  colback=GoldLight, colframe=Amber!40,
  fonttitle=\headingfont\bfseries, coltitle=GoldDark,
  title={#1},
  left=4mm, right=3.5mm, top=2.5mm, bottom=2.5mm
}

% Worked example box
\newtcolorbox{workedexbox}[1]{
  breakable, arc=2.5mm, boxrule=0.5pt,
  colback=white, colframe=Amber!50,
  fonttitle=\headingfont\bfseries, coltitle=GoldDark,
  title={#1},
  left=4mm, right=3.5mm, top=2.5mm, bottom=2.5mm
}
% Worked examples: \begin{workedex}{title} -> "Worked Example N.M --- title",
% auto-numbered per section and labelled we:N.M. Never type the number.
\newcounter{wenum}[section]
\renewcommand{\thewenum}{\thesection.\arabic{wenum}}
\newenvironment{workedex}[1]
  {\refstepcounter{wenum}\label{we:\thewenum}\begin{workedexbox}{Worked Example \thewenum\ --- #1}}
  {\end{workedexbox}}

% We do: straight after each worked example, a near-identical question the
% class works through together, with its own writing space. Numbered to
% match the worked example above it and labelled wd:N.M.
%   \begin{wedo}{title} question ... \answerbox{h} \end{wedo}
\newtcolorbox{wedobox}[1]{
  enhanced, breakable=false, arc=2.5mm, boxrule=0.5pt,
  colback=white, colframe=TealGreen!55,
  fonttitle=\headingfont\bfseries, coltitle=TealGreen!75!black,
  title={#1},
  left=4mm, right=3.5mm, top=2.5mm, bottom=2.5mm
}
\newenvironment{wedo}[1]
  {\par\needspace{9\baselineskip}\begin{wedobox}{We do \thewenum\ --- #1}\label{wd:\thewenum}\ignorespaces}
  {\end{wedobox}}

% Year 7 box (blue)
\newtcolorbox{yr7box}{
  breakable, arc=2.5mm, boxrule=0.5pt,
  colback=LightBlue, colframe=Blue!40,
  left=4mm, right=3.5mm, top=2.5mm, bottom=2.5mm
}

% Year 8 extension box (violet)
\newtcolorbox{yr8box}{
  breakable, arc=2.5mm, boxrule=0.5pt,
  colback=PurpleLight, colframe=Violet!40,
  left=4mm, right=3.5mm, top=2.5mm, bottom=2.5mm
}

% Answer lines box
\newcommand{\answerbox}[1]{%
  \begin{tcolorbox}[
    enhanced, breakable, arc=2mm,
    colback=Paper, colframe=Line,
    left=2mm, right=2mm, top=1mm, bottom=1mm,
    boxrule=1pt,
    height=#1
  ]
  \end{tcolorbox}%
}

% Experiment header
\newtcolorbox{expheader}[3]{
  breakable=false, arc=2.5mm,
  colback=RowAlt, colframe=Amber!50, boxrule=0.5pt,
  fonttitle=\headingfont\bfseries\large, coltitle=NavyBlue,
  title={Experiment #1: #2},
  left=4mm, right=3.5mm, top=2.5mm, bottom=2.5mm,
  after upper={\par\smallskip
    {\headingfont\bfseries\color{GoldDark}Aim:}\ \textit{#3}},
}

% Reading comprehension passage box (green)
\newtcolorbox{readingbox}[1]{
  breakable, arc=2.5mm, boxrule=0.5pt,
  colback=white, colframe=Teal!50,
  fonttitle=\headingfont\bfseries, coltitle=TealGreen,
  title={#1},
  left=4mm, right=3.5mm, top=2.5mm, bottom=2.5mm,
  fontupper=\normalsize
}

% Skill-builder box (for beginner scaffolding)
\newtcolorbox{skillbox}[1]{
  breakable, arc=2.5mm, boxrule=0.5pt,
  colback=RowAlt, colframe=Blue!40,
  fonttitle=\headingfont\bfseries, coltitle=MidBlue,
  title={#1},
  left=4mm, right=3.5mm, top=2.5mm, bottom=2.5mm
}

% ── Custom Commands ───────────────────────────────────────────────────────────

% Question: \question{marks}{text}  (marks can be empty)
% Numbered automatically as Q<section>.<n>, reset every section.
% Do NOT write numbers by hand.
\newcounter{qnum}[section]
\renewcommand{\theqnum}{\thesection.\arabic{qnum}}
\newcommand{\question}[2]{%
  \refstepcounter{qnum}%
  \needspace{3\baselineskip}%
  \medskip
  \noindent\label{q:\theqnum}\textbf{Q\theqnum}\ifx&#1&\else\ \textcolor{DarkGrey}{(#1 mark\ifnum#1=1\else s\fi)}\fi\quad #2%
  \medskip
}

% Inline marks note
\newcommand{\markscount}[1]{\hfill\textcolor{DarkGrey}{\textit{(#1 mark\ifnum#1=1\else s\fi)}}}

% ── Practice set counter and question macros ─────────────────────────────────
\newcounter{pqnum}
\newcommand{\practiceset}[1]{%
  \needspace{6\baselineskip}
  \setcounter{pqnum}{0}%
  \vspace{6pt}
  {\headingfont\bfseries\large\color{GoldDark} Practice Set --- #1}\\[-2pt]
  {\color{Amber}\rule{\linewidth}{1pt}}\\[2pt]
}
% \pq{question text}{working box height}
\newcommand{\pq}[2]{%
  \stepcounter{pqnum}%
  \needspace{4\baselineskip}
  \vspace{5pt}
  \noindent\textbf{\thepqnum.}\quad #1\par
  \vspace{2pt}
  \answerbox{#2}
}
% Two short questions side by side, each with its own box
\newcommand{\pqpair}[4]{%
  \needspace{4\baselineskip}
  \vspace{5pt}
  \noindent
  \begin{minipage}[t]{0.48\linewidth}
    \stepcounter{pqnum}\textbf{\thepqnum.}\quad #1\par\vspace{2pt}
    \begin{tcolorbox}[enhanced,arc=2mm,colback=Paper,colframe=Line,left=2mm,right=2mm,
      top=1mm,bottom=1mm,boxrule=1pt,height=#2]\end{tcolorbox}
  \end{minipage}\hfill
  \begin{minipage}[t]{0.48\linewidth}
    \stepcounter{pqnum}\textbf{\thepqnum.}\quad #3\par\vspace{2pt}
    \begin{tcolorbox}[enhanced,arc=2mm,colback=Paper,colframe=Line,left=2mm,right=2mm,
      top=1mm,bottom=1mm,boxrule=1pt,height=#4]\end{tcolorbox}
  \end{minipage}
}

% ── Comprehension question macro (compact, lined) ────────────────────────────
\newcounter{rcnum}
\newcommand{\rcreset}{\setcounter{rcnum}{0}}
\newcommand{\rcq}[2]{%
  \stepcounter{rcnum}%
  \needspace{3\baselineskip}
  \vspace{4pt}
  \noindent\textbf{C\thercnum.}\quad #1\par
  \vspace{2pt}
  \linedbox{#2}
}

% Section divider (big numbered banner): \sectiondiv{title}{summary}
%
% This IS the section: it steps the real section counter, so subsections,
% tasks, questions and figures inside it are numbered N.M automatically.
% Put \label{sec:N} straight after it. Never hand-type a section number.
\newcommand{\sectiondiv}[2]{%
  \needspace{5\baselineskip}%
  \refstepcounter{section}%
  % PDF bookmark for the lesson; its anchor, lesson.N.1, is what the contents table links to
  \pdfbookmark[1]{Lesson \thesection{} — #1}{lesson.\thesection}%
  \begin{tcolorbox}[
    enhanced, breakable=false, arc=3.5mm,
    colback=NavyBlue, colframe=NavyBlue, boxrule=0pt,
    sidebyside, sidebyside align=center,
    lefthand width=2.4cm,
    left=2mm, right=4mm, top=4mm, bottom=4mm
  ]
    \begin{center}
      {\headingfont\bfseries\color{white!35!NavyBlue}\fontsize{50}{50}\selectfont \thesection}
    \end{center}
    \tcblower
    {\headingfont\bfseries\color{white}\LARGE #1}\\[3pt]
    {\itshape\color{white!75!NavyBlue}\small #2}
  \end{tcolorbox}%
}

% Tasks: \task, \task[a] (part a of a NEW task), \taskpart{b} (next part
% of the same task), \exttask (Year 8 "Extension" — same counter).
% Numbered Task <section>.<n>, reset every section. Do NOT type numbers.
% Every task, question and figure also gets an automatic label
% (task:1.7b, q:2.4, fig:1.2, fig:E2.1) so page refs come from the .aux.
\newcounter{tasknum}[section]
\renewcommand{\thetasknum}{\thesection.\arabic{tasknum}}
\NewDocumentCommand{\task}{o}{\refstepcounter{tasknum}\label{task:\thetasknum\IfValueT{#1}{#1}}\textbf{Task \thetasknum\IfValueT{#1}{#1}}}
\newcommand{\taskpart}[1]{\label{task:\thetasknum#1}\textbf{Task \thetasknum#1}}
\newcommand{\exttask}{\refstepcounter{tasknum}\label{task:\thetasknum}\textbf{Extension \thetasknum}}

% Figure captions: \figcap[fig:label]{caption} -> "Figure N.M --- caption"
% \expfigcap{caption} -> "Figure EN.M" (experiment pages). Both auto-numbered.
\newcounter{fignum}[section]
\renewcommand{\thefignum}{\thesection.\arabic{fignum}}
\newcounter{expfignum}[section]
\renewcommand{\theexpfignum}{E\thesection.\arabic{expfignum}}
\newcommand{\figcap}[2][]{%
  \begin{center}
    \refstepcounter{fignum}\label{fig:\thefignum}\if\relax\detokenize{#1}\relax\else\label{#1}\fi
    \small\itshape\color{DarkGrey} Figure \thefignum\ --- #2
  \end{center}%
}
\newcommand{\expfigcap}[1]{%
  \begin{center}
    \refstepcounter{expfignum}\label{fig:\theexpfignum}%
    \small\itshape\color{DarkGrey} Figure \theexpfignum\ --- #1
  \end{center}%
}

% Diagram include
\newcommand{\diagramfig}[3]{%  {filename}{width fraction}{caption}
  \begin{center}
    \includegraphics[width=#2\linewidth]{#1}\\[4pt]
    \small\itshape\color{DarkGrey} #3
  \end{center}%
}

% Table column types
\newcolumntype{L}[1]{>{\RaggedRight\arraybackslash}p{#1}}
\newcolumntype{C}[1]{>{\Centering\arraybackslash}p{#1}}

% Table header colour
\newcommand{\thead}[1]{\cellcolor{TableHead}\color{white}\bfseries #1}
\newcommand{\talt}{\rowcolor{RowAlt}}
\newcommand{\tmid}{\rowcolor{MidGrey}}

% Inline year tags
\newcommand{\tagyr}[1]{%
  \tikz[baseline=-0.5ex]\node[
    rounded corners=2pt,
    fill=#1!20, draw=#1, text=#1,
    inner xsep=4pt, inner ysep=1pt,
    font=\scriptsize\bfseries
  ]{Yr #1 only};}

% Draw answer lines
\newcommand{\anslines}[1]{%
  \vspace{2pt}%
  \foreach \n in {1,...,#1}{%
    \par\noindent\rule{\linewidth}{0.3pt}\vspace{3pt}%
  }%
}

% ── Document ─────────────────────────────────────────────────────────────────
% ── Optional-lesson note, shown under the lesson banner ──────────────────────
\newcommand{\optionalnote}[1]{%
  % the title is a bold lead-in on the first line rather than a title bar,
  % which keeps Lesson 11's reading on the lesson's first page
  \begin{tcolorbox}[enhanced, breakable=false, arc=2.5mm, boxrule=0.5pt,
    colback=PurpleLight, colframe=Violet!50,
    left=4mm, right=3.5mm, top=2mm, bottom=2mm]
  {\headingfont\bfseries\color{PurpleDark}Optional lesson \textbullet{} Year 8 extension.}\enspace #1
  \end{tcolorbox}\vspace{2pt}}


% ══════════════════════════════════════════════════════════════════════════════
%  LESSON STRUCTURE (lesson booklet, Oct 2026)
%  Every lesson runs in the same four parts, in this order:
%    \lessonpart{A}{Reading and comprehension}{Teal}   \label{read:N}
%    \lessonpart{B}{Theory and examples}{Blue}         \label{theory:N}
%    \lessonpart{C}{Practice}{Amber}                    \label{prac:N}
%    \lessonpart{D}{Above level}{Violet}                \label{ext:N}
%  Above-level / higher-order questions live ONLY in Part D, at the end.
% ══════════════════════════════════════════════════════════════════════════════
\newcommand{\lessonpart}[3]{%
  \par\if#1A\needspace{10\baselineskip}\vspace{10pt}\else
    \if#1C\notesfill\fi\clearpage\fi
  \begin{tcolorbox}[enhanced, breakable=false, arc=2mm, boxrule=0pt,
    colback=#3!14, borderline west={4pt}{0pt}{#3},
    left=4mm, right=3mm, top=1.6mm, bottom=1.6mm]
  {\headingfont\bfseries\color{#3!60!black}\large Part #1 \textbullet{} #2}
  \end{tcolorbox}\vspace{2pt}}

% Key vocabulary at the start of each lesson: students copy each definition
% from the vocabulary slide in the lesson's starter. First row = red scaffold.
\newcommand{\vocabhead}{%
  \par\vspace{6pt}%
  \begin{tcolorbox}[enhanced, breakable=false, arc=2mm, boxrule=0pt,
    colback=NavyBlue!10, borderline west={4pt}{0pt}{NavyBlue},
    left=4mm, right=3mm, top=1.6mm, bottom=1.6mm]
  {\headingfont\bfseries\color{NavyBlue}\large Key vocabulary}
  \end{tcolorbox}\vspace{2pt}%
  \noindent Copy the definition of each word from the slide. The first one has been done for you.\par}

% Curriculum codes under a lesson banner
\newcommand{\lessoncodes}[1]{%
  {\small\headingfont\color{DarkGrey}Victorian Curriculum 2.0:\ \textbf{#1}}\par\vspace{4pt}}

% Two short tasks side by side, each with its own working box
% \taskpair{text}{box height}{text}{box height}
% The two boxes line up: the shorter question is padded by the difference in
% text height (measured in a scratch box, task counter restored, no label).
\newsavebox\tpbox
\newlength\tpA \newlength\tpB
\newcount\tpsaved
\newcommand\tpmeasure[2]{%
  \tpsaved=\value{tasknum}%
  \sbox\tpbox{\vbox{\hsize=0.485\linewidth \linewidth=\hsize \parskip=0pt
    \renewcommand\label[1]{}\task #1\par}}%
  \setlength{#2}{\dimexpr\ht\tpbox+\dp\tpbox\relax}%
  \setcounter{tasknum}{\tpsaved}}
\newcommand{\taskpair}[4]{%
  \par\needspace{9\baselineskip}\vspace{4pt}%
  \tpmeasure{#1}{\tpA}\tpmeasure{#3}{\tpB}%
  \noindent
  \edef\abgrp{\number\numexpr\value{abox}+1\relax}%
  \begin{minipage}[t]{0.485\linewidth}
    \task #1\par\ifdim\tpB>\tpA\vspace{\dimexpr\tpB-\tpA\relax}\fi\vspace{2pt}\answerbox{#2}
  \end{minipage}\hfill
  \begin{minipage}[t]{0.485\linewidth}
    \task #3\par\ifdim\tpA>\tpB\vspace{\dimexpr\tpA-\tpB\relax}\fi\vspace{2pt}\answerbox{#4}
  \end{minipage}\par\gdef\abgrp{0}\vspace{4pt}}
% Two parts of one question side by side, with their answer boxes level:
% \ptpair{a}{part (a) text}{box height}{b}{part (b) text}{box height}.
% The shorter part's text gets space below it to match the longer one.
\newcommand\ptmeasure[3]{%
  \sbox\tpbox{\vbox{\hsize=0.485\linewidth \linewidth=\hsize \parskip=0pt \pt{#1}#2\par}}%
  \setlength{#3}{\dimexpr\ht\tpbox+\dp\tpbox\relax}}
\newcommand{\ptpair}[6]{%
  \ptmeasure{#1}{#2}{\tpA}\ptmeasure{#4}{#5}{\tpB}%
  \par\noindent
  \edef\abgrp{\number\numexpr\value{abox}+1\relax}%
  \begin{minipage}[t]{0.485\linewidth}
    \pt{#1}#2\par\ifdim\tpB>\tpA\vspace{\dimexpr\tpB-\tpA\relax}\fi\answerbox{#3}
  \end{minipage}\hfill
  \begin{minipage}[t]{0.485\linewidth}
    \pt{#4}#5\par\ifdim\tpA>\tpB\vspace{\dimexpr\tpA-\tpB\relax}\fi\answerbox{#6}
  \end{minipage}\par\gdef\abgrp{0}}

% Answer boxes never split across a page (a half box on its own page is useless)
% ── Page filling ────────────────────────────────────────────────────────────
% Every writing box gets an id. Its page is logged to \jobname.abp; fill.py
% measures the empty space left at the foot of each page and shares it out
% among that page's boxes via \jobname.abx (\abset{id}{extra}). Side-by-side
% boxes share a group so they grow together.
\newcounter{abox}
\newwrite\abpos \immediate\openout\abpos=\jobname.abp
\newcommand{\abset}[2]{\expandafter\gdef\csname abx@#1\endcsname{#2}}
\InputIfFileExists{\jobname.abx}{}{}
\def\abgrp{0}
\newcommand{\abnext}{\stepcounter{abox}%
  \edef\abcur{\ifcsname abx@\theabox\endcsname\csname abx@\theabox\endcsname\else 0pt\fi}}
\newcommand{\ablog}{\edef\abtmp{\noexpand\write\abpos{\theabox\space\abgrp\space\noexpand\thepage}}\abtmp}
\newcommand{\abh}[1]{\dimexpr#1+\abcur\relax}
\renewcommand{\answerbox}[1]{%
  \abnext
  \begin{tcolorbox}[enhanced, breakable=false, arc=2mm, colback=Paper, colframe=Line,
    left=2mm, right=2mm, top=1mm, bottom=1mm, boxrule=1pt, height=\abh{#1}]
  \end{tcolorbox}\ablog%
}
% \linedbox{height}: an answer box ruled in dotted thirds. Lines are \thirdrow/3 apart and run dotted, dotted,
% solid (the baseline), repeating, and end on one dotted line. The number of
% rows is worked out from the box's final height, so the ruling still fills the
% box after fill.py has grown it; the rows are centred in the box.
\newlength\thirdrow \setlength\thirdrow{9mm}
\tikzset{thirdsolid/.style={draw=DarkGrey!70, line width=0.5pt},
         thirddot/.style={draw=DarkGrey!70, line width=0.6pt, dash pattern=on 0.9pt off 1.6pt}}
\newlength\lbH \newlength\lbO \newcount\lbN
\newcommand{\linedbox}[1]{%
  \abnext
  % interior height = box height less the 1pt frame top and bottom
  \setlength\lbH{\dimexpr\abh{#1}-2pt\relax}%
  % whole rows that fit with at least 3 mm spare top and bottom (at least one)
  \lbN=\numexpr(\dimexpr\lbH-6mm-\thirdrow/2\relax)/\thirdrow\relax
  \ifnum\lbN<1 \lbN=1 \fi
  \setlength\lbO{\dimexpr(\lbH-\lbN\thirdrow)/2\relax}%
  \begin{tcolorbox}[enhanced, breakable=false, arc=2mm, colback=Paper, colframe=Line,
    left=2mm, right=2mm, top=1mm, bottom=1mm, boxrule=1pt, height=\abh{#1},
    underlay={\begin{tcbclipinterior}
      \foreach \i in {1,...,\the\lbN}{
        \draw[thirddot] ([xshift=3mm,yshift=-\lbO-(\i-1)*\thirdrow]interior.north west) -- ([xshift=-3mm,yshift=-\lbO-(\i-1)*\thirdrow]interior.north east);
        \draw[thirddot] ([xshift=3mm,yshift=-\lbO-(\i-1)*\thirdrow-\thirdrow/3]interior.north west) -- ([xshift=-3mm,yshift=-\lbO-(\i-1)*\thirdrow-\thirdrow/3]interior.north east);
        \draw[thirdsolid] ([xshift=3mm,yshift=-\lbO-(\i-1)*\thirdrow-2*\thirdrow/3]interior.north west) -- ([xshift=-3mm,yshift=-\lbO-(\i-1)*\thirdrow-2*\thirdrow/3]interior.north east);}
      \draw[thirddot] ([xshift=3mm,yshift=-\lbO-\lbN\thirdrow]interior.north west) -- ([xshift=-3mm,yshift=-\lbO-\lbN\thirdrow]interior.north east);
    \end{tcbclipinterior}}]
  \end{tcolorbox}\ablog%
}
% Every written-response box is ruled the same way (Part A comprehension, We do,
% Part C tasks and Part D extensions, including side-by-side pairs). Space for
% mathematical working uses \workbox{height} instead and stays plain, as do
% \drawbox and the write-in tables.
\let\workbox\answerbox   % plain box: mathematical working (calculations) stays unruled
\let\answerbox\linedbox
% \drawbox{height}{contents}: white box for a diagram or drawing
\newcommand{\drawbox}[2]{%
  \abnext
  \begin{tcolorbox}[colback=white, colframe=MidGrey, boxrule=0.5pt, height=\abh{#1}]#2\end{tcolorbox}\ablog}
% Notes box at the end of Part B: fills whatever is left of the page (if worth it)
\newcommand{\notesfill}{\par
  \ifdim\dimexpr\pagegoal-\pagetotal\relax>3.2cm
    \vspace*{\fill}%
    \begin{tcolorbox}[enhanced, breakable=false, arc=2mm, colback=Paper, colframe=Line,
      left=2mm, right=2mm, top=1mm, bottom=1mm, boxrule=1pt,
      height=\dimexpr\pagegoal-\pagetotal-1.0cm\relax]
      {\small\headingfont\color{DarkGrey}Notes}
    \end{tcolorbox}%
  \fi}



% Centred table: wrap every content table in \begin{ctab} ... \end{ctab}
% (figures already centre themselves inside their \diag... macros)
% ctab also registers write-in tables (S column) for page filling: the extra
% height is per row, logged with the row count.
\newenvironment{ctab}{\par\smallskip\begingroup\centering
  \abnext\let\tbcur\abcur\global\wtrows0 }{\par
  \ifnum\wtrows>0 \edef\abtmp{\noexpand\write\abpos{\theabox\space\abgrp\space\noexpand\thepage\space\the\wtrows}}\abtmp\fi
  \endgroup\smallskip}

% ── Booklet formatting rules (Oct 2026) ─────────────────────────────────────
% 1. No marks are shown anywhere in the booklet.
\renewcommand{\markscount}[1]{}
% 2. A task's heading sits on its own line; its text starts on the next
%    line, left aligned. Parts (a), (b)… each start a new line: use \pt{a}.
\RenewDocumentCommand{\task}{o}{\par\noindent\refstepcounter{tasknum}\label{task:\thetasknum\IfValueT{#1}{#1}}\textbf{Task \thetasknum\IfValueT{#1}{#1}}\par\nobreak\noindent\ignorespaces}
\renewcommand{\taskpart}[1]{\par\noindent\label{task:\thetasknum#1}\textbf{Task \thetasknum#1}\par\nobreak\noindent\ignorespaces}
\renewcommand{\exttask}{\par\noindent\refstepcounter{tasknum}\label{task:\thetasknum}\textbf{Extension \thetasknum}\par\nobreak\noindent\ignorespaces}
\newcommand{\pt}[1]{\par\noindent\textbf{(#1)}\ \ignorespaces}
% Workstation pages (Experiment 1): numbered method steps with a hanging
% indent, and a one-line safety reminder under the materials list.
\newcommand{\wsstep}[1]{\par\noindent\hangindent=1.6em\hangafter=1\makebox[1.6em][l]{\textbf{#1}}\ignorespaces}
\newcommand{\wssafety}[1]{\par\smallskip{\small\textcolor{ScaffoldRed}{\headingfont\bfseries Safety:}\ #1}\par}
% 3. Tables have vertical rules between columns and down both sides
%    (write | between and around the columns, e.g. {|L{3cm}|C{2cm}|}).
%    Keep the normal padding at the outer edges (no @{} there): colortbl paints
%    header colour one \tabcolsep past each edge, so trimming the padding makes
%    the header overhang the grid.
\arrayrulecolor{MidGrey}
\setlength{\arrayrulewidth}{0.6pt}
% 4. Write-in tables: W/Z are left/centred columns that centre text
%    vertically; put the S column first (see below).
\newcolumntype{W}[1]{>{\RaggedRight\arraybackslash}m{#1}}
\newcolumntype{Z}[1]{>{\Centering\arraybackslash}m{#1}}
\def\tbcur{0pt}
\newcount\wtrows
\newcommand{\wh}{\global\advance\wtrows1 \rule[\dimexpr-0.75cm-\tbcur/2\relax]{0pt}{\dimexpr1.5cm+\tbcur\relax}}
% S = zero-width first column that makes every row of a write-in table
% at least 1.5 cm tall with its text vertically centred. Start each data
% row with "&"; start the header row with "\hd &" so it stays slim.
\newif\ifnowh
\newcolumntype{S}{@{}>{}c<{\ifnowh\else\wh\fi\global\nowhfalse}@{\hspace{\tabcolsep}}}
\newcommand{\hd}{\global\nowhtrue}
% 5. The first row of a multi-row write-in table is a worked scaffold in red.
\definecolor{ScaffoldRed}{HTML}{C0392B}
\newcommand{\scaf}[1]{\textcolor{ScaffoldRed}{\itshape #1}}
% 6. Tables are a light grid (rules between rows and columns) on white —
%    row shading is off, because colortbl paints shading over column rules.
\renewcommand{\talt}{}
\renewcommand{\tmid}{}
% ══════════════════════════════════════════════════════════════════════════════
%  DIAGRAM DEFINITIONS — The Particle Model
%  Self-contained TikZ macros, no external image files. Each \diag... macro is
%  called from the body by name.
% ══════════════════════════════════════════════════════════════════════════════
\usetikzlibrary{arrows.meta,calc,positioning,patterns,backgrounds}
\tikzset{
  thin arrow/.style={-{Stealth[length=6pt,width=4pt]}, line width=1.4pt},
}

% ── Lesson 1: the same three jars holding a solid, a liquid and a gas ───────────
% The solid keeps its own shape, the liquid takes the jar's shape with a flat top,
% the gas fills the whole jar. No particles yet: those arrive in Lesson 4.
\newcommand{\jar}[1]{% jar outline with its base at (#1,0)
  \draw[line width=1.2pt, Ink] (#1-1.25,2.9) -- (#1-1.25,0.15) arc (180:270:0.15)
    -- (#1+1.1,0) arc (270:360:0.15) -- (#1+1.25,2.9);}
\newcommand{\diagStates}{%
\begin{center}
\begin{tikzpicture}[font=\small, x=1cm, y=1cm]
  % solid: a block sitting in the corner
  \fill[GoldLight] (-0.95,0.02) rectangle (0.15,1.25);
  \draw[GoldDark, line width=1.1pt] (-0.95,0.02) rectangle (0.15,1.25);
  \jar{0}
  \node[font=\bfseries, text=NavyBlue] at (0,-0.45) {Solid};
  \node[align=center, font=\footnotesize, text=DarkGrey] at (0,-1.05) {keeps its own shape\\and volume};
  % liquid: takes the jar's shape, flat top
  \begin{scope}[xshift=4.2cm]
    \fill[Sky!35] (-1.25,0.15) arc (180:270:0.15) -- (1.1,0) arc (270:360:0.15) -- (1.25,1.1) -- (-1.25,1.1) -- cycle;
    \draw[MidBlue, line width=1.1pt] (-1.25,1.1) -- (1.25,1.1);
    \jar{0}
    \node[font=\bfseries, text=NavyBlue] at (0,-0.45) {Liquid};
    \node[align=center, font=\footnotesize, text=DarkGrey] at (0,-1.05) {takes the jar's shape,\\keeps its volume};
  \end{scope}
  % gas: fills the whole (sealed) jar
  \begin{scope}[xshift=8.4cm]
    \fill[Violet!16] (-1.25,0.15) arc (180:270:0.15) -- (1.1,0) arc (270:360:0.15) -- (1.25,2.9) -- (-1.25,2.9) -- cycle;
    \draw[line width=2.2pt, DarkGrey] (-1.4,2.9) -- (1.4,2.9);
    \jar{0}
    \node[font=\scriptsize, text=DarkGrey, anchor=south] at (0,2.95) {lid};
    \node[font=\bfseries, text=NavyBlue] at (0,-0.45) {Gas};
    \node[align=center, font=\footnotesize, text=DarkGrey] at (0,-1.05) {spreads out to fill\\the whole jar};
  \end{scope}
\end{tikzpicture}
\end{center}}

% ── Lesson 1, Worked Example 1.2 and We do 1.2: the same syringe, before and after
\newcommand{\diagSyringe}[3]{% {start mL}{end mL}{contents: gas|liquid}
\begin{center}\resizebox{\linewidth}{!}{%
\begin{tikzpicture}[font=\scriptsize, x=0.11cm, y=1cm]
  \foreach \row/\vol/\lab in {0/#1/before, -1.15/#2/after}{
    \begin{scope}[yshift=\row cm]
      \ifx#3g\fill[Violet!16] (0,0) rectangle (\vol,0.7);\else\fill[Sky!35] (0,0) rectangle (\vol,0.7);\fi
      \draw[line width=1pt, Ink] (0,0) rectangle (60,0.7);
      \draw[line width=1pt, Ink] (0,0.25) -- (-4,0.25) (0,0.45) -- (-4,0.45);
      \fill[Ink] (-5.5,0.15) rectangle (-4,0.55);
      \foreach \t in {0,10,...,60}{\draw[Ink] (\t,0.7) -- (\t,0.85); \node[anchor=south, font=\tiny] at (\t,0.82) {\t};}
      \fill[DarkGrey] (\vol,0.04) rectangle (\vol+1.6,0.66);
      \draw[line width=2pt, DarkGrey] (\vol+1.6,0.35) -- (\vol+26,0.35);
      \draw[line width=2pt, DarkGrey] (\vol+26,0.05) -- (\vol+26,0.65);
      \node[anchor=east, font=\scriptsize\bfseries, text=DarkGrey] at (-7,0.35) {\lab};
    \end{scope}}
  \node[font=\tiny, text=DarkGrey, anchor=west] at (61,0.35) {mL};
  \node[font=\tiny, text=DarkGrey, anchor=north] at (-4.75,-1.2) {sealed tip};
\end{tikzpicture}}
\end{center}}

% needspace.sty reserves room by planting a penalty -100 with 9 lines of
% temporary stretch. On a page that can't be filled tightly that fake break
% always wins, so a We do that fitted was pushed over and left half a page
% empty. This version breaks only when the room really isn't there, and
% ragged bottom lets every other page run on until something won't fit;
% fill.py then grows the writing boxes into what is left at the foot.
\renewcommand{\needspace}[1]{\par\ifdim\dimexpr\pagegoal-\pagetotal\relax<\dimexpr#1\relax\newpage\fi}
\raggedbottom
% Part D sits in a breakable yr8box, where \needspace can't see the page (the box
% is built off the page first). Wrap each extension in keepext so it is one
% unsplittable block: the yr8box then breaks between extensions, never inside one.
\newenvironment{keepext}{\par\begin{tcolorbox}[blankest, breakable=false]}{\end{tcolorbox}\par}
\usetikzlibrary{patterns}
\providecommand{\lessonart}{\thesection}

\providecommand{\unitname}{The Particle Model}\providecommand{\unitcaps}{THE PARTICLE MODEL}
% ══════════════════════════════════════════════════════════════════════════════
%  2027 WORKBOOK STYLE — A5 LESSON BOOKLET (added by make_a5.py)
%  Same look as the Space A5 trial and the A4 experiment booklets: colour
%  cover, black-and-white inside, Outfit headings, a lesson opener band with
%  the lesson's ink symbol, boxes told apart by frame and icon:
%    I do ● solid · We do ●● dashed · You do ●●● heavy · key idea / reading ◆
%    thin · Above level ▲ double · Country ◎ dotted
%  The body still names the master's colours; they are mapped to ink here.
% ══════════════════════════════════════════════════════════════════════════════
\usepackage{pdflscape}
\usetikzlibrary{patterns}

% A picture wider than the A5 text block is scaled down to fit it.
\newsavebox\fitb
\newenvironment{fitbox}{\begin{lrbox}{\fitb}}{\end{lrbox}%
  \ifdim\wd\fitb>\linewidth\resizebox{\linewidth}{!}{\usebox\fitb}\else\usebox\fitb\fi}

% ── Fonts: Outfit for headings ──────────────────────────────────────────────
\renewfontfamily\headingfont{Outfit}[
  Path=../fonts/, Extension=.ttf, Ligatures=TeX,
  UprightFont=Outfit-Regular, BoldFont=Outfit-SemiBold,
  ItalicFont=Outfit-Regular, ItalicFeatures={FakeSlant=0.18},
  BoldItalicFont=Outfit-SemiBold, BoldItalicFeatures={FakeSlant=0.18} ]

% ── Black and white inside ──────────────────────────────────────────────────
\definecolor{Ink}{HTML}{111111}
\colorlet{NavyBlue}{Ink}\colorlet{MidBlue}{Ink}\colorlet{TealGreen}{Ink}
\colorlet{GoldDark}{Ink}\colorlet{PurpleDark}{Ink}\colorlet{OrangeDark}{Ink}
\colorlet{TableHead}{Ink}
\colorlet{Blue}{Ink}\colorlet{Teal}{Ink}\colorlet{Amber}{Ink}\colorlet{Coral}{Ink}
\colorlet{Violet}{Ink}\colorlet{Sky}{Ink}
\colorlet{LightBlue}{white}\colorlet{LightGreen}{white}\colorlet{GoldLight}{white}
\colorlet{PurpleLight}{white}\colorlet{OrangeLight}{white}\colorlet{RowAlt}{white}
\colorlet{LightGrey}{white}\colorlet{Paper}{white}
\colorlet{DarkGrey}{black!72}\colorlet{MidGrey}{black!45}\colorlet{Line}{black!45}
\arrayrulecolor{black!60}
% Particle kinds stay distinguishable without colour: white, light grey,
% dotted and dark.
\tikzset{
  pA/.style={circle, fill=white, draw=Ink, line width=0.6pt, minimum size=3.2mm, inner sep=0pt},
  pB/.style={circle, fill=black!35, draw=Ink, line width=0.6pt, minimum size=3.2mm, inner sep=0pt},
  pC/.style={circle, preaction={fill=white}, pattern=dots, draw=Ink, line width=0.6pt, minimum size=3.2mm, inner sep=0pt},
  pD/.style={circle, fill=black!75, draw=Ink, line width=0.6pt, minimum size=3.2mm, inner sep=0pt}}

% ── Page style: mirrored, running head on the outside ───────────────────────
\fancypagestyle{fancy}{\fancyhf{}%
  \renewcommand{\headrulewidth}{0.4pt}%
  \renewcommand{\headrule}{\hbox to\headwidth{\color{Ink}\leaders\hrule height \headrulewidth\hfill}}%
  \fancyhead[LE]{\headingfont\footnotesize \unitname}%
  \fancyhead[RO]{\headingfont\footnotesize\leftmark}%
  \fancyfoot[LE,RO]{\headingfont\footnotesize\thepage}}
\fancypagestyle{plain}{\fancyhf{}\renewcommand{\headrulewidth}{0pt}%
  \fancyfoot[LE,RO]{\headingfont\footnotesize\thepage}}
\fancypagestyle{opener}{\fancyhf{}\renewcommand{\headrulewidth}{0pt}%
  \fancyfoot[LE,RO]{\headingfont\footnotesize\thepage}}
\pagestyle{fancy}\markboth{}{}
\raggedbottom
\setlength{\emergencystretch}{2em}

\titleformat{\subsection}{\headingfont\large\bfseries}{\thesubsection}{0.6em}{}
\titleformat{\subsubsection}{\headingfont\normalsize\bfseries}{}{0em}{}

% ── Lesson opener band (replaces the navy section banner) ────────────────────
\renewcommand{\sectiondiv}[2]{%
  \clearpage\thispagestyle{opener}%
  \refstepcounter{section}\markboth{\thesection\ \textbullet\ #1}{\thesection\ \textbullet\ #1}%
  \pdfbookmark[1]{Lesson \thesection{} — #1}{lesson.\thesection}\hypertarget{lesson.\thesection.1}{}%
  \noindent\begin{minipage}[b]{\dimexpr\linewidth-36mm\relax}
    {\headingfont\bfseries\scriptsize\addfontfeatures{LetterSpace=14}\unitcaps\ \textbullet\ LESSON}\par\vspace{1mm}
    {\headingfont\bfseries\fontsize{40}{40}\selectfont\thesection}\par\vspace{1.5mm}
    {\headingfont\bfseries\fontsize{15}{18}\selectfont\RaggedRight #1\par}\vspace{1mm}
    {\footnotesize\itshape\RaggedRight #2\par}
  \end{minipage}\hfill
  \begin{minipage}[b]{34mm}\includegraphics[width=34mm]{../art/l\lessonart-ink.pdf}\end{minipage}\par
  \vspace{2.5mm}\noindent\rule{\linewidth}{1.6pt}\par\vspace{1mm}}
\renewcommand{\lessoncodes}[1]{{\small\itshape VC 2.0: #1}\par\vspace{2pt}}

% Lesson parts: a heading and a rule; B, C and D start a new page.
\renewcommand{\lessonpart}[3]{%
  \par\if#1A\needspace{10\baselineskip}\vspace{8pt}\else
    \if#1C\notesfill\fi\clearpage\fi
  {\headingfont\bfseries\Large Part #1 \textbullet\ #2}\par\vspace{0.5mm}\noindent\rule{\linewidth}{0.6pt}\par\vspace{1mm}}
\renewcommand{\vocabhead}{%
  \par\vspace{4pt}{\headingfont\bfseries\Large Key vocabulary}\par\vspace{1pt}%
  \noindent{\small Copy the meaning of each word from the slide. The first one has been done for you.}\par}
\renewcommand{\thead}[1]{\headingfont\bfseries #1}

% ── Box system ──────────────────────────────────────────────────────────────
\newcommand{\icodots}[1]{\tikz[baseline=-0.6ex]{\foreach \k in {1,...,#1}{\fill[white] (\k*1.9mm,0) circle (0.7mm);}}}
\newcommand{\icokey}{\tikz[baseline=-0.6ex]{\fill[white] (0,0.9mm)--(0.9mm,0)--(0,-0.9mm)--(-0.9mm,0)--cycle;}}
\newcommand{\icoup}{\tikz[baseline=-0.6ex]{\fill[white] (-1mm,-0.8mm)--(1mm,-0.8mm)--(0,1mm)--cycle;}}
\newcommand{\icocountry}{\tikz[baseline=-0.6ex]{\draw[white,line width=0.5pt] (0,0) circle (1mm); \fill[white] (0,0) circle (0.45mm);}}
\newcommand{\icoread}{\tikz[baseline=-0.6ex]{\draw[white,line width=0.6pt] (-1.1mm,-0.8mm) rectangle (1.1mm,0.8mm); \draw[white,line width=0.5pt] (0,-0.8mm) -- (0,0.8mm);}}
\newcommand{\tagicon}[1]{\tikz[baseline=-0.6ex]\node[fill=Ink, rounded corners=0.8mm, inner xsep=1.2mm, inner ysep=1mm]{#1};}
\tcbset{wbbase/.style={enhanced, colback=white, arc=1.2mm,
    left=3mm, right=3mm, top=4.2mm, bottom=2.4mm, before skip=10pt, after skip=7pt,
    attach boxed title to top left={xshift=3mm, yshift=-2.4mm},
    boxed title style={colback=Ink, colframe=Ink, arc=0.8mm, boxrule=0pt,
                       left=1.6mm, right=1.8mm, top=0.6mm, bottom=0.6mm},
    fonttitle=\headingfont\bfseries\footnotesize, coltitle=white}}
\renewtcolorbox{workedexbox}[1]{wbbase, breakable, colframe=Ink, boxrule=1pt,
  title={\icodots{1}\ I do \textbullet\ #1}}
\renewenvironment{workedex}[1]
  {\refstepcounter{wenum}\label{we:\thewenum}\begin{workedexbox}{Worked example \thewenum\ --- #1}}
  {\end{workedexbox}}
\renewtcolorbox{wedobox}[1]{wbbase, breakable=false, frame hidden,
  borderline={1.1pt}{0pt}{Ink, dash pattern=on 4pt off 2.6pt},
  title={\icodots{2}\ #1}}
\renewenvironment{wedo}[1]
  {\par\begin{wedobox}{We do \thewenum\ --- #1}\label{wd:\thewenum}\ignorespaces}
  {\end{wedobox}}
\renewtcolorbox{yr7box}{wbbase, breakable, colframe=Ink, boxrule=2.4pt, top=5mm,
  title={\icodots{3}\ You do \textbullet\ Check your understanding}}
\renewtcolorbox{yr8box}{wbbase, breakable, frame hidden,
  borderline={0.6pt}{0pt}{Ink}, borderline={0.6pt}{1.6pt}{Ink},
  title={\icoup\ Above level \textbullet\ Year 8 extension}}
\renewtcolorbox{infobox}[1]{wbbase, breakable, colframe=Ink, boxrule=0.5pt, title={\icokey\ #1}}
\renewtcolorbox{scaffold}[1]{wbbase, breakable, colframe=Ink, boxrule=0.5pt, title={\icokey\ #1}}
\renewtcolorbox{skillbox}[1]{wbbase, breakable, colframe=Ink, boxrule=0.5pt, title={\icokey\ #1}}
\renewtcolorbox{readingbox}[1]{wbbase, breakable, colframe=Ink, boxrule=0.5pt, title={\icoread\ #1}}
\renewtcolorbox{expheader}[3]{wbbase, breakable=false, colframe=Ink, boxrule=1pt,
  title={\icokey\ Experiment #1: #2},
  after upper={\par\smallskip{\headingfont\bfseries Aim:}\ \textit{#3}}}
\newtcolorbox{countrybox}[1]{wbbase, breakable, frame hidden,
  borderline={1.4pt}{0pt}{Ink, line cap=round, dash pattern=on 0pt off 2.4pt},
  title={\icocountry\ #1}}

% ── Padding to a multiple of 4 (copier booklet mode) ────────────────────────
\newcommand{\notespage}{\clearpage\markboth{Notes}{Notes}%
  {\headingfont\bfseries\Large Notes}\par\vspace{3mm}%
  \noindent\begin{tikzpicture}\foreach \y in {0,...,17}{\draw[black!35, line width=0.4pt] (0,-\y*9mm) -- (\linewidth,-\y*9mm);}\end{tikzpicture}}
\newcount\padleft
\newcommand{\padtofour}{\clearpage
  \pgfmathtruncatemacro\pagesused{\the\value{page}-1}%
  \pgfmathtruncatemacro\padn{mod(4-mod(\pagesused,4),4)}%
  \padleft=\padn\relax
  \loop\ifnum\padleft>0 \notespage\advance\padleft by -1 \repeat
  \clearpage}

% ══════════════════════════════════════════════════════════════════════════════
%  WATER CYCLE AND RESOURCES — COVER
%  The unit's cover art (covers/water-resources.jpg), full bleed. The A5 lesson
%  booklet sets \wrcoverfile one folder deeper.
% ══════════════════════════════════════════════════════════════════════════════
\providecommand{\wrcoverfile}{../../covers/water-resources.jpg}
\newcommand{\wrcover}{%
\begin{tikzpicture}[remember picture, overlay]
  \node[anchor=north] at (current page.north) {\includegraphics[width=\paperwidth]{\wrcoverfile}};
\end{tikzpicture}}
% ── A5: every question full page width (stacked, never side by side) ───────
\renewcommand{\taskpair}[4]{%
  \par\needspace{6\baselineskip}\vspace{2pt}%
  \task #1\par\vspace{2pt}\answerbox{#2}%
  \par\needspace{6\baselineskip}%
  \task #3\par\vspace{2pt}\answerbox{#4}\par\vspace{2pt}}
\renewcommand{\ptpair}[6]{%
  \pt{#1}#2\par\answerbox{#3}%
  \pt{#4}#5\par\answerbox{#6}}
\renewcommand{\pqpair}[4]{%
  \needspace{4\baselineskip}\vspace{5pt}%
  \stepcounter{pqnum}\noindent\textbf{\thepqnum.}\quad #1\par\vspace{2pt}
  \begin{tcolorbox}[enhanced,arc=2mm,colback=Paper,colframe=Line,left=2mm,right=2mm,
    top=1mm,bottom=1mm,boxrule=1pt,height=#2]\end{tcolorbox}
  \stepcounter{pqnum}\noindent\textbf{\thepqnum.}\quad #3\par\vspace{2pt}
  \begin{tcolorbox}[enhanced,arc=2mm,colback=Paper,colframe=Line,left=2mm,right=2mm,
    top=1mm,bottom=1mm,boxrule=1pt,height=#4]\end{tcolorbox}}

\begin{document}
\pdfbookmark[1]{Cover}{cover}
\setlength{\parindent}{0pt}
\setlength{\parskip}{4pt}

% ---- Cover: the unit's cover art, full bleed, in colour ------------------
\thispagestyle{empty}
\wrcover
\begin{tikzpicture}[remember picture, overlay]
  
  \node[anchor=north east, fill=white, draw=Ink, line width=0.6pt, rounded corners=1.2mm,
        inner xsep=2mm, inner ysep=1.8mm, text width=30mm, font=\headingfont\footnotesize]
    at ([xshift=-4mm, yshift=-4mm]current page.north east)
    {\textbf{Name}\par\vspace{4mm}\hrule\vspace{2mm}\textbf{Class}\par\vspace{4mm}\hrule};
  \node[anchor=south east, fill=white, draw=Ink, line width=0.6pt, rounded corners=1.2mm,
        inner xsep=2.6mm, inner ysep=2mm, font=\headingfont\bfseries, align=right]
    at ([xshift=-4mm, yshift=6mm]current page.south east)
    {Lesson Booklet\\[-1pt]{\normalfont\headingfont\scriptsize Lessons 1, 3, 4 and 5}};
\end{tikzpicture}
\mbox{}
\clearpage

% ---- Inside cover ---------------------------------------------------------
\thispagestyle{plain}
{\headingfont\bfseries\fontsize{22}{26}\selectfont Water Cycle and Resources}\par\vspace{1mm}
{\headingfont\large Lesson Booklet \textbullet\ Year 7}\par\vspace{1mm}
{\small\itshape VC 2.0: VC2S8U09 \textbullet\ VC2S8U15 \textbullet\ VC2S8U16 \textbullet\ VC2S8H03}\par\vspace{4mm}
\textbf{Teacher}\ \hrulefill\par

\begin{countrybox}{Acknowledgement of Country}
This booklet was created on the unceded Country of the Wurundjeri Woi-wurrung people of the Kulin Nation. We acknowledge Elders past, present and emerging.
\end{countrybox}

\begin{infobox}{How this booklet works}
One section for each lesson. Each opens with its key vocabulary, then:\par\smallskip
\renewcommand{\arraystretch}{1.25}
\begin{tabular}{@{}p{24mm}p{70mm}@{}}
\textbf{Part A} & a reading with comprehension questions \\
\textbf{Part B} & the ideas, worked through \\
\textbf{Part C} & practice on your own \\
\end{tabular}

\smallskip
\begin{tabular}{@{}p{24mm}p{70mm}@{}}
\textbf{I do} \tagicon{\icodots{1}} & your teacher works an example (solid frame) \\
\textbf{We do} \tagicon{\icodots{2}} & you try one together (dashed frame) \\
\textbf{Reading} \tagicon{\icoread} & a passage to read (thin frame) \\
\end{tabular}

\smallskip
Lessons 2 and 6 are practicals: they are in the \textbf{Experiment Booklet}. Finished your lesson work and want to go further? Ask your teacher about \textbf{Advanced Water and Resources}.
\end{infobox}

\clearpage
\markboth{Contents}{Contents}
{\headingfont\bfseries\LARGE Contents}\par\vspace{1mm}\noindent\rule{\linewidth}{0.6pt}\par\vspace{2mm}
\renewcommand{\arraystretch}{1.45}
\begin{tabular}{|C{1.1cm}|L{6.0cm}|C{2.4cm}|}
  \hline\thead{Lesson} & \thead{Title} & \thead{Page} \\ \hline
  1 & \hyperlink{lesson.1.1}{What is the Water Cycle?} & \hyperlink{lesson.1.1}{p.\ \pageref{sec:1}} \\ \hline
  2 & \textit{Experiment 1 --- Water Cycle in a Bag} & Experiment Booklet \\ \hline
  3 & \hyperlink{lesson.3.1}{Resources} & \hyperlink{lesson.3.1}{p.\ \pageref{sec:3}} \\ \hline
  4 & \hyperlink{lesson.4.1}{Renewability and Sustainability} & \hyperlink{lesson.4.1}{p.\ \pageref{sec:4}} \\ \hline
  5 & \hyperlink{lesson.5.1}{Energy and Renewable Energy} & \hyperlink{lesson.5.1}{p.\ \pageref{sec:5}} \\ \hline
  6 & \textit{Experiment 2 --- Household Energy Audit} & Experiment Booklet \\ \hline
  --- & \hyperlink{glossary}{Glossary} & \hyperlink{glossary}{\pageref{sec:glossary}} \\ \hline
\end{tabular}
\renewcommand{\arraystretch}{1}
\clearpage
% ══════════════════════════════════════════════════════════════════════════════
%  LESSON 1 — WHAT IS THE WATER CYCLE?                                  (Week 1 · P2)
% ══════════════════════════════════════════════════════════════════════════════
\setcounter{section}{0}
\sectiondiv{What is the Water Cycle?}{Evaporation, condensation and precipitation \textbullet{} Collection and transpiration \textbullet{} Why fresh water is limited}\label{sec:1}
\lessoncodes{VC2S8U09 \textbullet{} VC2S8U05}

\vocabhead
\begin{ctab}
\begin{tabular}{|SW{3.4cm}|W{7.0cm}|}
  \hline\hd & \thead{Word} & \thead{Definition (copy from the slide)} \\ \hline
  & \scaf{\textbf{Water cycle}\label{voc:water-cycle}} & \scaf{The continuous movement of water between the land, the oceans and the air.} \\ \hline
  & \textbf{Evaporation}\label{voc:evaporation} & \\ \hline
  & \textbf{Condensation}\label{voc:condensation} & \\ \hline
  & \textbf{Precipitation}\label{voc:precipitation} & \\ \hline
  & \textbf{Transpiration}\label{voc:transpiration} & \\ \hline
  & \textbf{Renewable resource}\label{voc:renewable-resource} & \\ \hline
\end{tabular}
\end{ctab}

% ─── PART A ───────────────────────────────────────────────────────────────────
\lessonpart{A}{Reading and comprehension}{Teal}\label{read:1}

\begin{readingbox}{Reading --- The Same Water, Again and Again}
The water you drink today may once have fallen as rain on a dinosaur. Earth does not make new water; it recycles the same water over and over in the \textbf{water cycle}.

The Sun drives the cycle. Its energy warms the oceans, lakes and rivers, and water at the surface turns into water vapour, a gas. This is \textbf{evaporation}. Plants add more water vapour through tiny holes in their leaves: \textbf{transpiration}.

The warm, moist air rises. High in the sky it cools, and the water vapour turns back into tiny liquid droplets that gather into clouds. This is \textbf{condensation}. When the droplets join and grow too heavy, they fall as rain, hail or snow: \textbf{precipitation}.

Rain that lands on the ground runs into creeks and rivers, soaks into the soil to become groundwater, or collects in lakes, dams and the sea. Then the cycle starts again.

Because the cycle never stops, water is a \textbf{renewable resource}. Yet fresh water is still scarce. About 97\% of Earth's water is salty sea water, and most of the fresh water is frozen in ice sheets. Less than 1\% of all the water on Earth is fresh water we can easily reach.
\end{readingbox}

\rcreset
\rcq{What drives the water cycle?}{2.5cm}
\rcq{What is the difference between evaporation and transpiration?}{2.5cm}
\rcq{Explain how clouds form.}{2.5cm}
\rcq{Name three forms of precipitation.}{2.5cm}
\rcq{Give three places rain goes after it reaches the ground.}{2.5cm}
\rcq{If water is a renewable resource, why is fresh water still limited?}{3cm}

% ─── PART B ───────────────────────────────────────────────────────────────────
\lessonpart{B}{Theory and examples}{Blue}\label{theory:1}

\subsection{The Processes of the Water Cycle}\label{sub:1.1}

\begin{center}
\begin{fitbox}\begin{tikzpicture}[scale=0.92, font=\small]
  \fill[black!12] (6.2,0) rectangle (11,0.9);
  \draw[Ink, line width=0.8pt] (6.2,0.9) -- (11,0.9);
  \draw[Ink, line width=1pt, fill=black!5] (0,0) -- (2.4,3.0) -- (4.6,0.9) -- (6.2,0.9) -- (6.2,0) -- cycle;
  \draw[Ink, line width=1pt] (4.2,0.9) -- (4.2,1.9); \draw[Ink, line width=1pt, fill=white] (4.2,2.3) circle (0.45);
  \draw[Ink, line width=0.8pt, densely dashed] (1.4,0.6) .. controls (3,0.3) and (5,0.5) .. (6.4,0.5);
  \fill[Ink] (2.2,5.0) circle (0.5); \fill[Ink] (2.9,5.3) circle (0.6); \fill[Ink] (3.6,5.0) circle (0.45); \fill[Ink] (2.2,4.6) rectangle (3.6,5.0);
  \fill[Ink] (8.2,5.0) circle (0.45); \fill[Ink] (8.8,5.25) circle (0.55); \fill[Ink] (9.4,5.0) circle (0.4); \fill[Ink] (8.2,4.6) rectangle (9.4,5.0);
  \foreach \x in {2.3,2.8,3.3}{\draw[Ink, line width=1pt] (\x,4.3) -- (\x-0.2,3.6);}
  \draw[Ink, line width=1pt] (10.6,6.0) circle (0.45);
  \foreach \x in {7.4,8.6,9.8}{\draw[Ink, line width=1.2pt, -{Stealth[length=5pt]}, densely dotted] (\x,1.0) -- (\x,3.8);}
  \draw[Ink, line width=1.2pt, -{Stealth[length=5pt]}, densely dotted] (4.2,2.8) -- (4.6,4.3);
  \draw[Ink, line width=1.2pt, -{Stealth[length=5pt]}] (7.9,5.0) -- (4.0,5.1);
  \node[draw=Ink, circle, inner sep=1.5pt, fill=white, font=\footnotesize\bfseries] at (8.9,2.4) {1};
  \node[draw=Ink, circle, inner sep=1.5pt, fill=white, font=\footnotesize\bfseries] at (4.95,3.6) {2};
  \node[draw=Ink, circle, inner sep=1.5pt, fill=white, font=\footnotesize\bfseries] at (8.8,4.3) {3};
  \node[draw=Ink, circle, inner sep=1.5pt, fill=white, font=\footnotesize\bfseries] at (1.95,3.9) {4};
  \node[draw=Ink, circle, inner sep=1.5pt, fill=white, font=\footnotesize\bfseries] at (5.2,1.35) {5};
  \node[draw=Ink, circle, inner sep=1.5pt, fill=white, font=\footnotesize\bfseries] at (3.2,0.35) {6};
  \node[font=\scriptsize] at (8.6,0.45) {sea};
  \node[font=\scriptsize] at (10.6,6.65) {Sun};
\end{tikzpicture}\end{fitbox}
\figcap{The water cycle. Each number marks one process.}
\end{center}

\begin{skillbox}{The processes}
\textbf{Evaporation:} liquid water turns into water vapour, using energy from the Sun. \quad \textbf{Transpiration:} plants release water vapour from their leaves. \quad \textbf{Condensation:} water vapour cools and turns into liquid droplets (clouds). \quad \textbf{Precipitation:} water falls from clouds. \quad \textbf{Collection:} water runs off into rivers and the sea. \quad \textbf{Infiltration:} water soaks into the ground (groundwater).
\end{skillbox}

\begin{workedex}{Processes as changes of state}
Which change of state happens in evaporation and in condensation?

\smallskip
\textbf{Evaporation:} liquid $\rightarrow$ gas (water $\rightarrow$ water vapour). It needs energy, from the Sun.\\
\textbf{Condensation:} gas $\rightarrow$ liquid (water vapour $\rightarrow$ droplets). It happens when the vapour cools.
\end{workedex}

\begin{wedo}{Processes as changes of state}
Snow falls on a mountain, and in spring it melts into a creek. Name the process in the water cycle and the change of state for each step.
\answerbox{4cm}
\end{wedo}

\subsection{Fresh Water is Limited}\label{sub:1.2}

\begin{workedex}{Earth's water as a bucket}
Imagine all of Earth's water is a 10-litre bucket. About 97\% is salty: $0.97 \times 10 = 9.7$ litres. The fresh water is only about 0.3 litres, and most of that is frozen. The fresh water we can easily reach is less than one cupful.
\end{workedex}

\begin{wedo}{Earth's water as a bucket}
Use the bucket to explain to a friend why we should not waste fresh water, even though the water cycle keeps renewing it.
\answerbox{4cm}
\end{wedo}

% ─── PART C ───────────────────────────────────────────────────────────────────
\lessonpart{C}{Practice}{Amber}\label{prac:1}

\task Name the process at each number on a water cycle diagram. The processes are: evaporation from the sea (1), transpiration from a tree (2), condensation into a cloud (3), precipitation from a cloud (4), collection as water runs off into the sea (5), and infiltration into the ground (6). Write the process for each number, and the change of state if there is one.
\begin{ctab}
\begin{tabular}{|SZ{1.2cm}|W{3.8cm}|W{4.6cm}|}
  \hline\hd & \thead{No.} & \thead{Process} & \thead{Change of state} \\ \hline
  & \scaf{1} & \scaf{Evaporation} & \scaf{Liquid to gas} \\ \hline
  & 2 & & \\ \hline
  & 3 & & \\ \hline
  & 4 & & \\ \hline
  & 5 & & \\ \hline
  & 6 & & \\ \hline
\end{tabular}
\end{ctab}

\needspace{10\baselineskip}
\task Wet washing dries faster on a hot, windy day. Explain why, using the word \emph{evaporation}.
\answerbox{5cm}

\needspace{10\baselineskip}
\task A cold can of drink gets wet on the outside on a warm day. Explain where the water comes from, using the word \emph{condensation}.
\answerbox{5cm}

\needspace{10\baselineskip}
\task Explain why water is called a renewable resource, but fresh water is still limited.
\answerbox{5cm}
\newpage


% ══════════════════════════════════════════════════════════════════════════════
%  LESSON 3 — RESOURCES                                                 (Week 2 · P1)
% ══════════════════════════════════════════════════════════════════════════════
\setcounter{section}{2}
\sectiondiv{Resources}{Renewable and non-renewable resources \textbullet{} How resources are extracted \textbullet{} Benefits and risks of extraction}\label{sec:3}
\lessoncodes{VC2S8U09 \textbullet{} VC2S8H03}

\vocabhead
\begin{ctab}
\begin{tabular}{|SW{3.4cm}|W{7.0cm}|}
  \hline\hd & \thead{Word} & \thead{Definition (copy from the slide)} \\ \hline
  & \scaf{\textbf{Resource}\label{voc:resource}} & \scaf{Something from the natural world that people use, such as water, wood, metals or coal.} \\ \hline
  & \textbf{Renewable}\label{voc:renewable} & \\ \hline
  & \textbf{Non-renewable}\label{voc:non-renewable} & \\ \hline
  & \textbf{Ore}\label{voc:ore} & \\ \hline
  & \textbf{Fossil fuel}\label{voc:fossil-fuel} & \\ \hline
  & \textbf{Extraction}\label{voc:extraction} & \\ \hline
\end{tabular}
\end{ctab}

% ─── PART A ───────────────────────────────────────────────────────────────────
\lessonpart{A}{Reading and comprehension}{Teal}\label{read:3}

\begin{readingbox}{Reading --- Digging Up Australia}
Australia is one of the world's great mining countries. Iron ore from the Pilbara in Western Australia is shipped to steelworks around the world. Coal is mined in Queensland, New South Wales and in Victoria's Latrobe Valley, and natural gas is piped from under the sea. These are all \textbf{resources}: materials from the natural world that people use.

Some resources are \textbf{renewable}: they are replaced by nature in a human lifetime, like timber from a well-managed forest, or energy from the Sun and wind. Others are \textbf{non-renewable}. Coal, oil and natural gas are \textbf{fossil fuels}, formed over millions of years from the remains of ancient plants and animals. Once they are used, they are gone.

Getting a resource out of the ground is called \textbf{extraction}. Metals are found in rocks called \textbf{ores}, which are blasted and dug from mines, then crushed and processed to separate the metal. Extraction brings great benefits: jobs, money for the country, and the materials for everything from phones to bridges. It also brings risks: cleared land, lost habitat, polluted water and dust, and greenhouse gases.
\end{readingbox}

\rcreset
\rcq{Name three resources mined in Australia, and where each is found.}{2.5cm}
\rcq{What makes a resource renewable?}{2.5cm}
\rcq{Why are fossil fuels non-renewable?}{2.5cm}
\rcq{What is an ore?}{2.5cm}
\rcq{Give two benefits of extraction.}{2.5cm}
\rcq{Give two risks of extraction.}{2.5cm}

% ─── PART B ───────────────────────────────────────────────────────────────────
\lessonpart{B}{Theory and examples}{Blue}\label{theory:3}

\subsection{Renewable or Non-Renewable?}\label{sub:3.1}

\begin{skillbox}{The test}
Ask: \emph{will nature replace it within a human lifetime?} If yes, it is \textbf{renewable}. If it takes thousands or millions of years, it is \textbf{non-renewable}. A renewable resource can still run out if it is used faster than it is replaced.
\end{skillbox}

\begin{workedex}{Classifying resources}
\textbf{Timber:} renewable --- trees regrow in decades, if forests are replanted.\\
\textbf{Gold:} non-renewable --- the Earth makes no new gold in a human lifetime.\\
\textbf{Fish:} renewable, but only if not caught faster than they breed.
\end{workedex}

\begin{wedo}{Classifying resources}
Classify sunlight and natural gas as renewable or non-renewable, and give a reason for each.
\answerbox{4cm}
\end{wedo}

\subsection{Benefits and Risks of Extraction}\label{sub:3.2}

\begin{workedex}{An open-cut coal mine}
\textbf{Benefits:} jobs for the town; coal to generate electricity; money from selling coal.\\
\textbf{Risks:} cleared land and lost habitat; dust; carbon dioxide released when the coal is burned; a large hole to repair when the mine closes.
\end{workedex}

\begin{wedo}{A new lithium mine}
Lithium is used in the batteries of phones and electric cars. Give one benefit and one risk of opening a new lithium mine.
\answerbox{4cm}
\end{wedo}

\begin{countrybox}{Country and extraction}
In 2020, mining in the Pilbara destroyed the Juukan Gorge rock shelters, sacred sites of the Puutu Kunti Kurrama and Pinikura peoples with evidence of people living there for more than 46\,000 years. Deciding whether to extract a resource means weighing benefits and risks for everyone affected, including the Traditional Owners of the land.
\end{countrybox}

% ─── PART C ───────────────────────────────────────────────────────────────────
\lessonpart{C}{Practice}{Amber}\label{prac:3}

\task Sort each resource as renewable (R) or non-renewable (N), and give a reason.
\begin{ctab}
\begin{tabular}{|SW{2.8cm}|Z{1.2cm}|W{5.8cm}|}
  \hline\hd & \thead{Resource} & \thead{R or N} & \thead{Reason} \\ \hline
  & \scaf{Coal} & \scaf{N} & \scaf{Takes millions of years to form.} \\ \hline
  & Wind & & \\ \hline
  & Iron ore & & \\ \hline
  & Cotton & & \\ \hline
  & Oil & & \\ \hline
  & Fresh water & & \\ \hline
\end{tabular}
\end{ctab}

\needspace{10\baselineskip}
\task Describe the steps to get a metal from its ore, from the mine to the metal.
\answerbox{5cm}

\needspace{12\baselineskip}
\task A company wants to open a gold mine near a country town.
\pt{a}Give two benefits for the town.
\answerbox{4cm}
\pt{b}Give two risks for the environment.
\answerbox{4cm}

\needspace{10\baselineskip}
\task Fish are a renewable resource, yet some fish populations have collapsed. Explain how this can happen.
\answerbox{5cm}
\newpage


% ══════════════════════════════════════════════════════════════════════════════
%  LESSON 4 — RENEWABILITY AND SUSTAINABILITY                           (Week 2 · P2)
% ══════════════════════════════════════════════════════════════════════════════
\setcounter{section}{3}
\sectiondiv{Renewability and Sustainability}{What makes a use sustainable \textbullet{} Reduce, reuse, recycle, substitute \textbullet{} Melbourne's water}\label{sec:4}
\lessoncodes{VC2S8U09 \textbullet{} VC2S8H03}

\vocabhead
\begin{ctab}
\begin{tabular}{|SW{3.4cm}|W{7.0cm}|}
  \hline\hd & \thead{Word} & \thead{Definition (copy from the slide)} \\ \hline
  & \scaf{\textbf{Sustainable}\label{voc:sustainable}} & \scaf{Using a resource in a way that leaves enough for people in the future.} \\ \hline
  & \textbf{Conserve}\label{voc:conserve} & \\ \hline
  & \textbf{Reduce}\label{voc:reduce} & \\ \hline
  & \textbf{Reuse}\label{voc:reuse} & \\ \hline
  & \textbf{Recycle}\label{voc:recycle} & \\ \hline
  & \textbf{Substitute}\label{voc:substitute} & \\ \hline
\end{tabular}
\end{ctab}

% ─── PART A ───────────────────────────────────────────────────────────────────
\lessonpart{A}{Reading and comprehension}{Teal}\label{read:4}

\begin{readingbox}{Reading --- When Melbourne Ran Low on Water}
Between about 1997 and 2009, south-eastern Australia went through the Millennium Drought, one of the longest dry spells on record. Less rain fell, and the rivers that fill Melbourne's dams slowed. By 2009 the dams had fallen to about a quarter full.

Melbourne had to use water more \textbf{sustainably}: in a way that would leave enough for the future. Households were asked to \textbf{conserve} water and cut their use to 155 litres per person each day. People \textbf{reduced} their use with shorter showers and water-saving taps. They \textbf{reused} water from the laundry on gardens. Some water was \textbf{recycled}, cleaned and used again on farms and parks. Rainwater tanks \textbf{substituted} for drinking water in toilets and gardens.

The state also built a desalination plant at Wonthaggi, which takes the salt out of sea water. It gives Melbourne a source of water that does not depend on rain, but it uses a great deal of energy to run.
\end{readingbox}

\rcreset
\rcq{What was the Millennium Drought?}{2.5cm}
\rcq{How full were Melbourne's dams by 2009?}{2.5cm}
\rcq{What daily target were households given?}{2.5cm}
\rcq{Give one example each of reducing, reusing and substituting water.}{3cm}
\rcq{What does a desalination plant do?}{2.5cm}
\rcq{Give one benefit and one drawback of the desalination plant.}{2.5cm}

% ─── PART B ───────────────────────────────────────────────────────────────────
\lessonpart{B}{Theory and examples}{Blue}\label{theory:4}

\subsection{What Makes a Use Sustainable?}\label{sub:4.1}

\begin{skillbox}{The sustainability test}
A use of a resource is \textbf{sustainable} if it can go on for a very long time without running the resource out or damaging the environment. For a renewable resource, use it no faster than it is replaced. For a non-renewable resource, use as little as possible and recover it where you can.
\end{skillbox}

\begin{workedex}{Is it sustainable?}
A town takes water from a river for its farms. In a dry year, it takes more water than flows into the river.

\smallskip
\textbf{Not sustainable:} it uses the water faster than the rain replaces it, so the river runs low and fish and plants downstream suffer.
\end{workedex}

\begin{wedo}{Is it sustainable?}
A school prints every worksheet single-sided and throws used paper in the general rubbish. Decide whether this is sustainable, and explain.
\answerbox{4cm}
\end{wedo}

\subsection{Four Ways to Conserve Resources}\label{sub:4.2}

\begin{skillbox}{Reduce, reuse, recycle, substitute}
\textbf{Reduce:} use less in the first place. \quad \textbf{Reuse:} use something again for the same or a new job. \quad \textbf{Recycle:} break a material down and make it into something new. \quad \textbf{Substitute:} use a different, more sustainable resource instead.
\end{skillbox}

\begin{workedex}{Plastic bottles}
\textbf{Reduce:} carry a refillable bottle. \textbf{Reuse:} refill a bottle you already have. \textbf{Recycle:} put empty bottles in the recycling bin to be made into new plastic. \textbf{Substitute:} use a steel bottle that lasts for years.
\end{workedex}

\begin{wedo}{School lunch packaging}
Give one way to reduce, reuse, recycle and substitute the packaging in a school lunch.
\answerbox{4.5cm}
\end{wedo}

% ─── PART C ───────────────────────────────────────────────────────────────────
\lessonpart{C}{Practice}{Amber}\label{prac:4}

\task Write whether each action is reduce, reuse, recycle or substitute.
\begin{ctab}
\begin{tabular}{|SW{6.8cm}|W{3.0cm}|}
  \hline\hd & \thead{Action} & \thead{Which one?} \\ \hline
  & \scaf{Turning off the tap while brushing teeth} & \scaf{Reduce} \\ \hline
  & Using a glass jar to store pencils & \\ \hline
  & Putting old newspaper in the recycling bin & \\ \hline
  & Using rainwater to flush the toilet & \\ \hline
  & Taking shorter showers & \\ \hline
\end{tabular}
\end{ctab}

\needspace{10\baselineskip}
\task Suggest three ways our school could use water more sustainably. Explain how each one helps.
\answerbox{5cm}

\needspace{10\baselineskip}
\task Is it sustainable to burn coal for electricity? Explain your answer.
\answerbox{5cm}

\needspace{10\baselineskip}
\task Explain why recycling is better than throwing things away, but reducing is better than recycling.
\answerbox{5cm}
\newpage


% ══════════════════════════════════════════════════════════════════════════════
%  LESSON 5 — ENERGY AND RENEWABLE ENERGY                               (Week 2 · P3)
% ══════════════════════════════════════════════════════════════════════════════
\setcounter{section}{4}
\sectiondiv{Energy and Renewable Energy}{How electricity is generated \textbullet{} Fossil fuels and renewables \textbullet{} Comparing energy sources}\label{sec:5}
\lessoncodes{VC2S8U09 \textbullet{} VC2S8U15 \textbullet{} VC2S8H03}

\vocabhead
\begin{ctab}
\begin{tabular}{|SW{3.4cm}|W{7.0cm}|}
  \hline\hd & \thead{Word} & \thead{Definition (copy from the slide)} \\ \hline
  & \scaf{\textbf{Generator}\label{voc:generator}} & \scaf{A machine that turns the spinning of a turbine into electricity.} \\ \hline
  & \textbf{Turbine}\label{voc:turbine} & \\ \hline
  & \textbf{Renewable energy}\label{voc:renewable-energy} & \\ \hline
  & \textbf{Greenhouse gas}\label{voc:greenhouse-gas} & \\ \hline
  & \textbf{Reliability}\label{voc:reliability} & \\ \hline
\end{tabular}
\end{ctab}

% ─── PART A ───────────────────────────────────────────────────────────────────
\lessonpart{A}{Reading and comprehension}{Teal}\label{read:5}

\begin{readingbox}{Reading --- Where Victoria's Electricity Comes From}
For most of the last century, Victoria's electricity came from brown coal dug from huge open-cut mines in the Latrobe Valley. In a coal power station, the coal is burned to boil water. The steam spins a \textbf{turbine}, and the turbine turns a \textbf{generator}, which makes electricity. Burning coal releases carbon dioxide, a \textbf{greenhouse gas} that traps heat in the atmosphere and adds to climate change.

Today more and more of Victoria's electricity comes from \textbf{renewable energy}. Wind turbines spin in the wind on hills and along the coast. Solar panels on roofs turn sunlight straight into electricity. Hydroelectric power stations in the mountains let water fall through turbines.

Each source has trade-offs. A coal power station can run day and night: its \textbf{reliability} is high. The Sun does not shine at night and the wind does not always blow, so renewable electricity has to be stored in batteries or balanced with other sources. Engineers now plan an electricity system that is clean, reliable and affordable all at once.
\end{readingbox}

\rcreset
\rcq{Where has most of Victoria's electricity come from in the past?}{2.5cm}
\rcq{Describe how a coal power station makes electricity.}{3cm}
\rcq{Why is burning coal a problem for the climate?}{2.5cm}
\rcq{Name three renewable sources of electricity.}{2.5cm}
\rcq{What does reliability mean for an energy source?}{2.5cm}
\rcq{Why does renewable electricity sometimes need to be stored?}{2.5cm}

% ─── PART B ───────────────────────────────────────────────────────────────────
\lessonpart{B}{Theory and examples}{Blue}\label{theory:5}

\subsection{How Electricity is Generated}\label{sub:5.1}

\begin{center}
\begin{fitbox}\begin{tikzpicture}[font=\small, box/.style={draw=Ink, line width=0.8pt, rounded corners=1.5mm, minimum width=17mm, minimum height=9mm, align=center, fill=white}]
  \node[box] (f) at (0,0) {Fuel\\burned};
  \node[box] (w) at (2.25,0) {Water\\boiled};
  \node[box] (s) at (4.5,0) {Steam\\spins turbine};
  \node[box] (g) at (6.75,0) {Generator};
  \node[box] (e) at (9.0,0) {Electricity};
  \foreach \a/\b in {f/w,w/s,s/g,g/e}{\draw[Ink, line width=1pt, -{Stealth[length=5pt]}] (\a) -- (\b);}
\end{tikzpicture}\end{fitbox}
\figcap{In a coal or gas power station. Wind and hydro spin the turbine directly; solar panels need no turbine at all.}
\end{center}

\begin{workedex}{Energy transfers in a wind turbine}
The wind's moving energy turns the blades (\textbf{turbine}). The spinning shaft turns the \textbf{generator}, which produces \textbf{electrical} energy. No fuel is burned, so no greenhouse gases are released while it runs.
\end{workedex}

\begin{wedo}{Energy transfers in a hydro power station}
Water stored behind a dam falls through pipes to the power station. Describe how this makes electricity.
\answerbox{4cm}
\end{wedo}

\subsection{Comparing Energy Sources}\label{sub:5.2}

\begin{workedex}{Coal and solar}
\textbf{Coal:} reliable day and night; existing power stations; releases large amounts of greenhouse gas; non-renewable.\\
\textbf{Solar:} renewable; no greenhouse gas while running; panels on homes cut power bills; makes no electricity at night, so needs storage or another source.
\end{workedex}

\begin{wedo}{Wind}
Give one benefit and one drawback of wind power.
\answerbox{4cm}
\end{wedo}

% ─── PART C ───────────────────────────────────────────────────────────────────
\lessonpart{C}{Practice}{Amber}\label{prac:5}

\task Compare four energy sources.
\begin{ctab}
\begin{tabular}{|SW{1.5cm}|Z{2.1cm}|W{3.0cm}|W{3.3cm}|}
  \hline\hd & \thead{Source} & \thead{Renewable?} & \thead{Benefit} & \thead{Risk or drawback} \\ \hline
  & \scaf{Coal} & \scaf{No} & \scaf{Reliable day and night} & \scaf{Releases greenhouse gases} \\ \hline
  & Solar & & & \\ \hline
  & Wind & & & \\ \hline
  & Hydro & & & \\ \hline
  & Gas & & & \\ \hline
\end{tabular}
\end{ctab}

\needspace{10\baselineskip}
\task Explain why a solar power system often has a battery.
\answerbox{5cm}

\needspace{10\baselineskip}
\task A town must choose between a new gas power station and a wind farm with batteries. Recommend one, giving two reasons.
\answerbox{5cm}

\needspace{10\baselineskip}
\task Explain why electricity from a coal power station is not sustainable in the long term.
\answerbox{5cm}
\newpage

\clearpage\markboth{Glossary}{Glossary}\hypertarget{glossary}{}\pdfbookmark[1]{Glossary}{glossarybm}
{\headingfont\bfseries\LARGE Glossary}\label{sec:glossary}\par\vspace{1mm}\noindent\rule{\linewidth}{0.6pt}\par
{\small Every key word in this booklet and the page it is introduced on.}\par\smallskip
\renewcommand{\arraystretch}{1.3}
\begin{longtable}{|>{\RaggedRight}p{34mm}|p{9mm}|p{\dimexpr\linewidth-43mm-6\tabcolsep-4\arrayrulewidth-1pt\relax}|}
\hline \rule{0pt}{4mm}\headingfont\bfseries Word & \headingfont\bfseries Page & \headingfont\bfseries In my own words \\ \hline \endhead
  Condensation & \pageref{voc:condensation} & \\ \hline
  Conserve & \pageref{voc:conserve} & \\ \hline
  Evaporation & \pageref{voc:evaporation} & \\ \hline
  Extraction & \pageref{voc:extraction} & \\ \hline
  Fossil fuel & \pageref{voc:fossil-fuel} & \\ \hline
  Generator & \pageref{voc:generator} & \\ \hline
  Greenhouse gas & \pageref{voc:greenhouse-gas} & \\ \hline
  Non-renewable & \pageref{voc:non-renewable} & \\ \hline
  Ore & \pageref{voc:ore} & \\ \hline
  Precipitation & \pageref{voc:precipitation} & \\ \hline
  Recycle & \pageref{voc:recycle} & \\ \hline
  Reduce & \pageref{voc:reduce} & \\ \hline
  Reliability & \pageref{voc:reliability} & \\ \hline
  Renewable & \pageref{voc:renewable} & \\ \hline
  Renewable energy & \pageref{voc:renewable-energy} & \\ \hline
  Renewable resource & \pageref{voc:renewable-resource} & \\ \hline
  Resource & \pageref{voc:resource} & \\ \hline
  Reuse & \pageref{voc:reuse} & \\ \hline
  Substitute & \pageref{voc:substitute} & \\ \hline
  Sustainable & \pageref{voc:sustainable} & \\ \hline
  Transpiration & \pageref{voc:transpiration} & \\ \hline
  Turbine & \pageref{voc:turbine} & \\ \hline
  Water cycle & \pageref{voc:water-cycle} & \\ \hline
\end{longtable}
\padtofour
\end{document}
